% !TeX root = ../../../poisson-cubic.tex
\section{特殊なブロック制限とFermat型}\label{sec:pg-special}

本節では、正規形\eqref{eq:pg-normal}において典型的な「滑らかな平面3次曲線」の仮定が成り立たない、あるいは退化が極端に集中する2つの重要な特殊ケースを調べる。第一に、あるブロックの3変数制限 $f_s(x_s,0)$ が恒等的に零となる場合に必要な中心変数の個数を決定する。第二に、Fermat型3次超曲面において零点スキームの残余部分 $Z$ がどのように非被約な厚化として現れるかを明示する。

\subsection{3変数制限が恒等的に零となる場合}
前節では3変数制限 $f_s(x_s,0)$ が特異な平面3次曲線（節点型・尖点型）になる例を構成した。さらに極端な状況として、$f_s(x_s,0)$ が恒等的に零になることも、中心変数の個数 $m$ が十分大きければ起こり得る。

\begin{proposition}\label{prop:pg-zero-restriction}
正規形\eqref{eq:pg-normal}で与えられる3次超曲面 $X$ が滑らかであり、あるブロック $s$ について3変数制限 $f_s(x_s,0)$ が恒等的に零であるとする。このとき $m\geq3$ でなければならない。逆に、$m=3$ の場合にはこのような滑らかな3次超曲面が存在する。
\end{proposition}
\begin{proof}
まず必要条件 $m\geq3$ を示す。$f_s(x_s,0)\equiv 0$ ならば、対応する座標平面 $\PP_s^2=\{u=0,\ x_t=0\ (t\neq s)\}$ 上で $F$ の全てのブロック変数に関する偏微分が恒等的に消える。残る偏微分は中心変数に関する高々 $m$ 個の2次斉次式 $F_{u_\alpha}(x_s,0)$（$1\le\alpha\le m$）だけである。もし $m\leq2$ ならば、$\PP_s^2$ 上でこれら高々2個の2次式の共通零点（すなわち3変数多項式環で高さ高々2のイデアルの非零零点）が必ず存在する。そのような点では $F$ の全偏微分が消えるため、$X$ の滑らかさに反する。ゆえに $m\geq3$ である。

次に $m=3$ での存在を具体例で示す。9変数 $a_1,a_2,a_3,b_1,b_2,b_3,u_1,u_2,u_3$ の3次式
\begin{equation}\label{eq:pg-zero-restriction-example}
F=\sum_{i=1}^3(u_i a_i^2+u_i^3)+b_1^3+b_2^3+b_3^3
\end{equation}
を考える。全偏微分が同時に零になるとすると、まず $F_{a_i}=2u_ia_i=0$ および $F_{u_i}=a_i^2+3u_i^2=0$ から各 $i$ について $a_i=u_i=0$ が従い、さらに $F_{b_i}=3b_i^2=0$ から $b_i=0$ となる。従って共通零点はアフィン原点だけであり、$X=V(F)\subset\PP^8$ は滑らかである。これは $k=2$、$m=3$ の正規形であり、$a$-ブロックの $u=0$ 制限は恒等的に零となっている。
\end{proof}

なお、例\eqref{eq:pg-zero-restriction-example}の退化スキームは、3変数制限の滑らかさを仮定しない一般的な生成元表示（式\eqref{eq:pg-deg-generators}）によってそのまま記述される。

\subsection{Fermat型における非被約な零点スキーム}
命題\ref{prop:pg-zero-ci}では、一般の係数に対して零点スキームの残余部分 $Z=\Delta\cap X$ が滑らかになることを見た。対照的に、Fermat型3次超曲面では臨界点被覆 $\Delta$ の全ての分岐が中心線形空間上に集中し、$Z$ は一様な厚みを持つ非被約スキームとなる。

正規形の座標を用いて、Fermat型の3次式を
\begin{equation}\label{eq:pg-fermat}
F=\frac13\left(\sum_{\alpha=1}^m u_\alpha^3+
\sum_{s=1}^k\sum_{i=1}^3x_{si}^3\right)
\end{equation}
と書く。
\begin{proposition}\label{prop:pg-fermat-thickening}
Fermat型3次式\eqref{eq:pg-fermat}において、$m>0$ ならば臨界点スキーム $\Delta=V(x_{si}^2\mid s,i)$ は中心線形空間 $\PP_u^{m-1}=\{x_{si}=0\}$ を台とし、横断方向の長さ $8^k$ を持つ厚化である。さらに $m\geq2$ ならば、残余零点スキーム $Z=\Delta\cap X$ は $\PP_u^{m-1}$ 内のFermat超曲面 $V(\sum_{\alpha=1}^m u_\alpha^3)$ の同じ横断長さ $8^k$ の厚化となる。一方、$m=0$ の場合には $\Delta=\varnothing$ である。
\end{proposition}
\begin{proof}
式\eqref{eq:pg-fermat}のブロック変数に関する偏微分は $A_{s,i}=x_{si}^2$ である。従って $\Delta$ の斉次座標環は
\begin{equation}
\CC[u]\otimes_\CC\bigotimes_{s=1}^k
\CC[x_{s1},x_{s2},x_{s3}]/(x_{s1}^2,x_{s2}^2,x_{s3}^2)
\end{equation}
と同型になる。各ブロックに対応するArtin局所環 $\CC[x_{s1},x_{s2},x_{s3}]/(x_{s1}^2,x_{s2}^2,x_{s3}^2)$ の長さは $2^3=8$ であるため、$k$ 個のテンソル積は横断長さ $8^k$ の厚化を与える。さらに、この商環では全ての $x_{si}^3$ が零になるため、$F$ の像は $\frac13\sum_{\alpha=1}^m u_\alpha^3$ に一致する。これが $m\ge2$ における $Z=\Delta\cap X$ の厚化構造を与える。最後に $m=0$ のとき、偏微分 $x_{si}^2=0$ の共通零点はアフィン原点のみであるため、射影スキーム $\Delta$ は空である。
\end{proof}
